% TEMPLATE-MILC-v3.tex - plantilla maestra UNICA de las guias MILC v3.
%
% La usan las cuatro familias de guias, cada una con su driver:
%   · Grados 6-11          -> scripts/build-guias-g11.py
%   · Bebras               -> scripts/build-guias-bebras.py
%   · Semillero            -> scripts/build-guias-semillero.py
%   · Territorio interior  -> scripts/build-guias-territorio-interior.py
%
% Antes habia tres copias casi identicas (~400 lineas iguales, 6 distintas):
% cada mejora habia que aplicarla tres veces, y en la practica no se hacia
% (el motor de recursos vivia solo en dos de ellas). Lo que varia entre
% familias esta parametrizado en 6 placeholders que cada driver llena
% (se nombran aqui SIN su sintaxis real: el builder reemplaza por texto plano,
% tambien dentro de los comentarios, y un valor multilinea rompe el preambulo):
%
%   HEADER_IZQ        encabezado izquierdo de pagina
%   HEADER_DER        encabezado derecho de pagina
%   PORTADA_ETIQUETA  rotulo sobre el numero grande (GRADO/MOMENTO/MODULO)
%   PORTADA_SERIE     linea de serie de la portada
%   PORTADA_ID        identificador de la guia en la portada
%   PORTADA_CIRCULO   el circulito de grado (°), o vacio si no aplica
%   PDF_TITULO        titulo en texto plano (sin \textbf ni comillas TeX) para
%                     los metadatos del PDF (/Title); lo produce lib_xelatex.texto_pdf
%
% Composicion (auditoria 2026-09-03): las bandas de estacion piden espacio
% (\Needspace*) y prohiben el salto justo despues (\nopagebreak), asi no quedan
% huerfanas al pie; las cajas largas (softbox, drawbox, infoband, puentebox) son
% `breakable`, asi una caja de media pagina ya no arrastra todo a la siguiente
% dejando la anterior al 40 %. Todo texto sobre mostaza o verde va en milcNegro
% (blanco sobre #E5B400 da 1,93:1 y sobre #58A923 2,95:1; ninguno pasa WCAG AA).
% El resto de claves las documenta content/guias/_SCHEMA.md. El script valida
% que no queden placeholders sin reemplazar antes de escribir el archivo final.

% Etiquetado PDF (latex-lab): /Lang, estructura y texto alternativo de las
% figuras, para lectores de pantalla. Debe ir ANTES de \documentclass.
\DocumentMetadata{lang=es-CO,tagging=on}
\documentclass[11pt,letterpaper]{article}
% headheight=24pt: la cabecera derecha lleva dos lineas (identidad de la guia
% y estacion actual). headsep se acorta en la misma medida para que el bloque
% de texto quede exactamente donde estaba con headheight=12pt.
\usepackage[margin=2.0cm,headheight=24pt,headsep=13pt]{geometry}
\usepackage[table]{xcolor}
\usepackage{fontspec}
\usepackage[spanish]{babel}
\usepackage{tikz}
% Librerías TikZ para los diagramas de las guías (bloque `recursos:`):
%  · babel  → babel en español vuelve activos caracteres como «>», y TikZ los usa
%             en sus opciones (>=stealth, ->). Sin esta librería, cualquier
%             diagrama con flechas revienta con «\language@active@arg> has an
%             extra }». Debe ir ANTES de que se dibuje nada.
%  · positioning        → sintaxis «below left=2pt and 3pt».
%  · decorations.pathreplacing → llaves ({brace}) para señalar rangos.
\usetikzlibrary{babel,positioning,decorations.pathreplacing}
\usepackage{graphicx}
% Circuitos: el trabajo de electrónica se hace hoy con micro:bit y sus sensores
% integrados (MakeCode, sin cableado), así que no se necesitan esquemáticos.
% Si algún día entran montajes con protoboard, basta descomentar esta línea:
% los diagramas .tex/.tikz declarados en `recursos:` ya se incrustan solos.
% \usepackage{circuitikz}
\usepackage[most]{tcolorbox}
\usepackage{tabularx}
\usepackage{array}
\usepackage{fancyhdr}
\usepackage{titlesec}
\usepackage[document]{ragged2e}  % bandera a la derecha en todo el cuerpo
\usepackage{enumitem}
\usepackage{microtype}
\usepackage{etoolbox}
\usepackage{needspace}
% hyperref al final del preambulo: metadatos del PDF (titulo, autor, idioma,
% familia), marcadores por estacion (\pdfbookmark) y enlaces sin recuadro.
\usepackage[hidelinks,
  pdftitle={Ciberbullying — prevención y Línea 141},
  pdfauthor={PhD. Álvaro Cárdenas Orozco},
  pdfsubject={Tecnología e Informática · Grado 8 · Guía 26 · MILC v3},
  pdflang={es-CO},
  bookmarksopen=true,bookmarksnumbered=false]{hyperref}

% Helvetica Neue no trae la flecha U+2192, asi que cada «→» escrito literalmente
% en el YAML se compilaba a NADA, en silencio (462 flechas en 61 guias: rutas del
% tipo «paleta → tipografia → lockup» salian rotas). Mapeamos el caracter a su
% version matematica, que si tiene glifo. Asi el autor puede seguir escribiendo
% «→» comodamente en el YAML.
\usepackage{newunicodechar}
\newunicodechar{→}{\ensuremath{\rightarrow}}
% Mismo problema con los simbolos de marca y vinetas: el «✓» de «una palomita ✓»
% salia como cuadrito vacio en el PDF de todas las guias que lo usaban.
\usepackage{pifont}
\newunicodechar{✓}{\ding{51}}
\newunicodechar{✗}{\ding{55}}
\newunicodechar{✔}{\ding{52}}
\newunicodechar{•}{\textbullet}
\newunicodechar{≥}{\ensuremath{\geq}}
\newunicodechar{≤}{\ensuremath{\leq}}

\setmainfont{Helvetica Neue}
\setsansfont{Helvetica Neue}
\setlength{\parindent}{0pt}
\setlength{\parskip}{6pt}
% Sin particion de palabras: en 6.º-7.º una palabra cortada por guion al final
% de linea («Comuni-cacion») frena la lectura mucho mas que una linea corta.
\hyphenpenalty=10000
\exhyphenpenalty=10000
\pretolerance=2000
\tolerance=3000
\setlength{\emergencystretch}{3em}
\linespread{1.10}
\renewcommand{\arraystretch}{1.25}
\setlist{leftmargin=5mm,itemsep=1mm,topsep=1mm}
\newcolumntype{Y}{>{\RaggedRight\arraybackslash}X}

\definecolor{milcMagenta}{HTML}{E6007E}
\definecolor{milcVino}{HTML}{5A0038}
\definecolor{milcTurquesa}{HTML}{0093A5}
\definecolor{milcVerde}{HTML}{58A923}
\definecolor{milcMostaza}{HTML}{E5B400}
\definecolor{milcOcre}{HTML}{B86600}
\definecolor{milcNegro}{HTML}{241816}
\definecolor{milcGris}{HTML}{F7F7F4}
\definecolor{milcLinea}{HTML}{E8E2DD}
\definecolor{milcGradePrimary}{HTML}{FF6600}
\definecolor{milcGradeDark}{HTML}{9A3E00}
\definecolor{milcGradeSoft}{HTML}{FFE4D0}
\definecolor{milcGradeAccent}{HTML}{CFE7FF}
\definecolor{milcGradeText}{HTML}{FFFFFF}
\color{milcNegro}

\pagestyle{fancy}
\fancyhf{}
% Cabecera de dos lineas: identidad de la guia arriba y, a la derecha y en
% gris, la estacion en curso (\markright desde \milcestacion). Antes de la
% primera estacion la segunda linea va vacia; el \strut mantiene la altura
% para que la linea de arriba no baile entre paginas.
\lhead{\small Tecnología e Informática · Grado 8\\[-1pt]{\footnotesize\strut}}
\rhead{\small Guía 26 · MILC v3\\[-1pt]{\footnotesize\strut\color{milcNegro!65}\rightmark}}
\cfoot{\small\thepage}
\renewcommand{\headrulewidth}{0pt}

% Sobre mostaza (#E5B400) y verde (#58A923) el texto va en milcNegro; sobre el
% resto de la paleta, en blanco. Se usa dentro de fonttitle/fontupper, donde un
% \color posterior manda sobre coltitle.
\newcommand{\milctintasobre}[1]{%
  \ifboolexpr{test{\ifstrequal{#1}{milcMostaza}} or test{\ifstrequal{#1}{milcVerde}}}%
    {\color{milcNegro}}{\color{white}}}

\newtcolorbox{titlebox}[1]{
  enhanced,colback=#1,colframe=#1,arc=7mm,boxrule=0pt,
  left=7mm,right=7mm,top=3mm,bottom=3mm,
  before skip=8pt,after skip=8pt,
  fontupper=\sffamily\bfseries\Large\color{white}
}

\newtcolorbox{softbox}[3]{
  enhanced,breakable,pad at break=3mm,
  colback=#2,colframe=#1,borderline west={4pt}{0pt}{#1},
  arc=2mm,boxrule=.4pt,left=4mm,right=4mm,top=3mm,bottom=3mm,
  before skip=6pt,after skip=8pt,title={#3},coltitle=white,
  colbacktitle=#1,fonttitle=\sffamily\bfseries\milctintasobre{#1},fontupper=\RaggedRight,
  attach boxed title to top left={xshift=3mm,yshift=-2mm},
  boxed title style={arc=2mm, boxrule=0pt}
}

\newtcolorbox{drawbox}[2]{
  enhanced,breakable,pad at break=4mm,
  colback=white,colframe=#1,arc=4mm,boxrule=1pt,
  left=5mm,right=5mm,top=4mm,bottom=4mm,before skip=8pt,after skip=8pt,
  title={#2},coltitle=white,colbacktitle=#1,fonttitle=\sffamily\bfseries\milctintasobre{#1},
  fontupper=\RaggedRight,
  attach boxed title to top left={xshift=4mm,yshift=-2mm},
  boxed title style={arc=3mm, boxrule=0pt}
}

\newtcolorbox{infoband}[1]{
  enhanced,breakable,pad at break=2mm,colback=#1!8,colframe=#1!50,arc=2mm,boxrule=.5pt,
  left=4mm,right=4mm,top=2mm,bottom=2mm,before skip=4pt,after skip=4pt,
  fontupper=\small\RaggedRight
}

% Puente narrativo entre secciones (contrato editorial Conéctate).
% Da continuidad: "Acabas de X. Ahora vas a Y porque Z."
% Visualmente: banda izquierda en color del grado, texto en itálica pequeña.
\newtcolorbox{puentebox}{
  enhanced,breakable,pad at break=2mm,colback=milcGradeSoft,colframe=milcGradeDark,
  borderline west={3pt}{0pt}{milcGradeDark},
  arc=0mm,boxrule=0pt,left=5mm,right=4mm,top=2.5mm,bottom=2.5mm,
  before skip=4pt,after skip=8pt,
  fontupper=\itshape\small\color{milcNegro}\RaggedRight
}
% La flecha va en modo matematico a proposito: Helvetica Neue NO tiene el glifo
% U+2192, asi que \textrightarrow se compilaba a NADA («Missing character: There
% is no → in font Helvetica Neue») y el puente salia sin flecha. En modo
% matematico la flecha viene de la fuente matematica y si se dibuja.
\newcommand{\puente}[1]{%
  \begin{puentebox}%
    \textcolor{milcGradeDark}{$\rightarrow$}\hspace{2mm}#1%
  \end{puentebox}%
}

\newcommand{\linead}[1]{\noindent\color{milcLinea}\rule{#1}{0.5pt}\color{milcNegro}}
\newcommand{\respuesta}[1]{\par\vspace{2mm}\foreach \n in {1,...,#1}{\noindent\color{milcLinea}\rule{\linewidth}{0.5pt}\par\vspace{5mm}}\color{milcNegro}}
\newcommand{\checkbox}{\textcolor{milcGradeDark}{\fbox{\phantom{x}}}\hspace{5pt}}

% ─── Listas aireadas para las guías socioemocionales ────────────────────
% Componer las verificaciones como «\checkbox texto\\» deja el texto que salta
% de línea debajo del cuadrito y apelmaza el bloque. Estos entornos alinean la
% continuación y airean los ítems. Los usa Territorio Interior; cualquier
% familia puede hacerlo.
\newlist{checklista}{itemize}{1}
\setlist[checklista]{
  label=\textcolor{milcGradeDark}{\fbox{\phantom{x}}},
  leftmargin=9mm, labelsep=3.2mm, itemsep=2.4mm,
  topsep=2.5mm, parsep=0pt, partopsep=0pt, before=\RaggedRight
}
\newlist{pasoslista}{enumerate}{1}
\setlist[pasoslista]{
  label=\textcolor{milcGradeDark}{\sffamily\bfseries\arabic*.},
  leftmargin=8mm, labelsep=3mm, itemsep=2.8mm,
  topsep=1mm, parsep=0pt, partopsep=0pt, before=\RaggedRight
}

\newcommand{\phasecard}[4]{
\begin{tcolorbox}[enhanced,colback=#3,colframe=#2,arc=3mm,boxrule=.5pt,height=3.05cm,valign=center,halign=center,left=1.5mm,right=1.5mm,top=2mm,bottom=2mm]
{\sffamily\bfseries\fontsize{22}{24}\selectfont\textcolor{#2}{#1}}\par
{\sffamily\bfseries\small #4}
\end{tcolorbox}
}

% ── Piezas del rediseno 2026 (legibilidad 6.º-11.º) ───────────────────
% Banda de estacion: reemplaza el titlebox macizo. Titulo a la izquierda,
% dato practico (tiempo/modalidad) a la derecha, en una sola linea.
% Nunca huerfana: pide 9 lineas libres (o salta de pagina rellenando con
% \vfill) y prohibe el salto justo despues de la banda. Nueve y no seis:
% la caja partible que sigue necesita ~80pt (titulo adosado, relleno y dos
% lineas) para arrancar en la misma pagina; con menos, tcolorbox la manda
% entera a la siguiente y la banda se queda sola. El penalty 10000 va
% en `after`, ANTES del espacio posterior: un `after skip` (aunque sea 0pt)
% es un \vskip tras la caja, y ese glue es un punto de corte legal que un
% \nopagebreak escrito despues de \end{tcolorbox} ya no cierra. Cada banda
% deja marcador PDF y actualiza la cabecera derecha.
\newcounter{milcestacion}
\newcommand{\milcestacion}[2]{%
\Needspace*{9\baselineskip}%
\stepcounter{milcestacion}%
\pdfbookmark[1]{#2}{milcestacion\themilcestacion}%
\markright{#2}%
\begin{tcolorbox}[enhanced,colback=#1,colframe=#1,arc=2mm,boxrule=0pt,
  left=5mm,right=5mm,top=2.6mm,bottom=2.6mm,before skip=11pt,
  after={\par\nopagebreak\vskip7pt}]
{\sffamily\bfseries\fontsize{15}{18}\selectfont\milctintasobre{#1}#2}
\end{tcolorbox}}

% Tarjeta de paso numerada: sustituye las tablas «Pilar / Aplicacion», que
% eran formato de consulta docente, no de estudiante. El numero va en un
% circulo sobre el borde para que la secuencia se lea de un vistazo.
\newcommand{\milcpaso}[3]{%
\Needspace{4\baselineskip}%
\begin{tcolorbox}[enhanced,colback=white,colframe=milcLinea,arc=1.5mm,boxrule=.9pt,
  left=14mm,right=4mm,top=3mm,bottom=3mm,before skip=4pt,after skip=4pt,
  fontupper=\RaggedRight,
  overlay={\node[anchor=north west,circle,fill=#2,text=white,inner sep=0pt,
    minimum size=7.4mm,font=\sffamily\bfseries\small]
    at ([xshift=3.6mm,yshift=-3.2mm]frame.north west) {#1};}]
#3
\end{tcolorbox}}

% Casilla marcable de tamano util con lapiz (la anterior era \fbox{\phantom{x}},
% de alto variable segun el texto que la seguia).
\newcommand{\milcmarca}{\framebox[3.8mm]{\rule{0pt}{3.4mm}}}

% Fila marcable de lista suelta (checks de la estacion 1).
\newcommand{\milccheckrow}[1]{%
\par\vspace{1.8mm}\noindent
\makebox[6.4mm][l]{\raisebox{-.55mm}{\textcolor{milcNegro}{\milcmarca}}}%
\parbox[t]{\dimexpr\linewidth-6.4mm\relax}{\RaggedRight #1}\par}

% Celda marcable dentro de la rubrica (tabla de autoevaluacion). Misma
% sangria francesa, pero pensada para vivir dentro de una columna X.
\newcommand{\milcopcion}[1]{%
\makebox[5.2mm][l]{\raisebox{-.55mm}{\textcolor{milcNegro}{\milcmarca}}}%
\parbox[t]{\dimexpr\linewidth-5.2mm\relax}{\RaggedRight #1}}

% ── Recursos visuales de la guía ──────────────────────────────────────
% Declarados en el bloque `recursos:` del YAML y emitidos por el builder.
% \guiaFigura   → imagen rasterizada o PDF (foto, carta celeste, montaje).
% \guiaDiagrama → fuente TikZ/circuitikz (.tex) incrustada como vector:
%                 es la vía para circuitos y esquemáticos editables.
% [#1] = texto alternativo (alt) · #2 = ruta relativa al .tex · #3 = pie
% El alt llega al PDF etiquetado (/Alt) y lo lee el lector de pantalla.
\newcommand{\guiaFigura}[3][]{%
\begin{center}
\includegraphics[width=0.88\linewidth,height=0.38\textheight,keepaspectratio,alt={#1}]{#2}\\[1.5mm]
{\footnotesize\itshape\textcolor{milcNegro!75}{#3}}
\end{center}
}

\newcommand{\guiaDiagrama}[2]{%
\begin{center}
\input{#1}\\[1.5mm]
{\footnotesize\itshape\textcolor{milcNegro!75}{#2}}
\end{center}
}

\begin{document}
\thispagestyle{empty}

% ── Portada ───────────────────────────────────────────────────────────
% Antes: un rectangulo de color a pagina completa. Se veia bien en pantalla,
% pero en la fotocopiadora del colegio se comia el toner y salia gris sucio.
% Ahora la banda ocupa el tercio superior y el resto va en blanco.
\begin{tikzpicture}[remember picture,overlay]
  \path[fill=milcGradePrimary]
    (current page.north west) rectangle ([yshift=-10.6cm]current page.north east);

  \node[anchor=north west,text=milcGradeText] at ([xshift=2.0cm,yshift=-1.55cm]current page.north west)
    {\sffamily\bfseries\fontsize{12}{14}\selectfont GRADO};
  \node[anchor=north west,text=milcGradeText,opacity=.85] at ([xshift=2.0cm,yshift=-2.35cm]current page.north west)
    {\sffamily\bfseries\fontsize{8.5}{10}\selectfont SERIE GUÍAS · TECNOLOGÍA E INFORMÁTICA};
  \node[anchor=north east,text=milcGradeText,opacity=.85,align=right] at ([xshift=-2.0cm,yshift=-1.55cm]current page.north east)
    {\sffamily\bfseries\fontsize{9}{12}\selectfont I.E. Sor María Juliana\\Cartago, Valle del Cauca};

  \node[anchor=west,text=milcGradeText] (gradonum) at ([xshift=1.85cm,yshift=-7.1cm]current page.north west)
    {\sffamily\bfseries\fontsize{112}{112}\selectfont 8};
  % El circulito de grado (°) solo aplica a las guias de grado («11°»); en
  % Bebras y Semillero el numero grande es un momento/modulo, no un grado.
  % Se posiciona relativo al nodo `gradonum`, no en coordenadas absolutas.
  \draw[milcGradeText,line width=1.6pt]
    ([xshift=2.0mm,yshift=-4.5mm]gradonum.north east) circle (.15cm);

  \node[anchor=west,text=milcGradeText,text width=11.4cm,align=left] at ([xshift=1.5cm,yshift=0.5cm]gradonum.east)
    {
     {\sffamily\bfseries\fontsize{9}{11}\selectfont\MakeUppercase{Período 3 · Multimedia, narrativa y ciberseguridad}}\\[3mm]
     {\sffamily\bfseries\fontsize{21}{25}\selectfont\setlength{\baselineskip}{27pt}Ciberbullying ---\\[4pt]prevención\\[4pt]y Línea 141}\\[3.5mm]
     {\sffamily\bfseries\fontsize{15}{18}\selectfont Guía 26-8-TIC}
    };
\end{tikzpicture}

\vspace*{9.4cm}
\vfill

{\sffamily\bfseries\fontsize{8}{10}\selectfont\textcolor{milcGradeDark}{LO QUE TE LLEVAS HOY}}

\begin{tcolorbox}[enhanced,colback=white,colframe=milcNegro,arc=2mm,boxrule=1.6pt,
  left=5mm,right=5mm,top=4mm,bottom=4mm,before skip=3pt,after skip=12pt,
  fontupper=\RaggedRight\fontsize{11}{15.5}\selectfont]
Tres piezas de convivencia digital: (a) el mapa de las cuatro fases del ciberbullying, de las primeras señales a la denuncia, con dos ejemplos concretos por fase; (b) un plan personal de cinco pasos con nombres reales de las personas a las que acudirías; (c) un cartel para el aula con la Línea 141 destacada y los demás canales, listo para colgar. Cierra con una ficha de cinco líneas donde declaras un compromiso con la convivencia del grupo.
\end{tcolorbox}

\begin{tcolorbox}[enhanced,colback=milcGris,colframe=milcGris,arc=2mm,boxrule=0pt,
  left=5mm,right=5mm,top=3.5mm,bottom=3.5mm,after skip=12pt]
{\sffamily\bfseries\fontsize{8}{10}\selectfont\textcolor{milcNegro!55}{ESTA GUÍA ES DE}}\\[3mm]
\begin{tabularx}{\linewidth}{X p{3.2cm} p{3.2cm}}
\linead{\linewidth} & \linead{3.0cm} & \linead{3.0cm}\\[-1mm]
{\fontsize{8.5}{10}\selectfont\textcolor{milcNegro!55}{Nombre completo}} &
{\fontsize{8.5}{10}\selectfont\textcolor{milcNegro!55}{Grupo}} &
{\fontsize{8.5}{10}\selectfont\textcolor{milcNegro!55}{Fecha}}\\
\end{tabularx}
\end{tcolorbox}

\vfill

\noindent\textcolor{milcLinea}{\rule{\linewidth}{1pt}}\\[2mm]
{\sffamily\bfseries\fontsize{9.5}{12}\selectfont Autor: Dr. Álvaro Cárdenas Orozco}\\[1mm]
{\sffamily\fontsize{8.5}{11}\selectfont\textcolor{milcNegro!55}{Metodología MILC · apertura ancestral · pensamiento computacional · triángulo Dussel–estoicismo–Floridi}}

\newpage

\section*{Apertura: Ciberbullying --- prevención y Línea 141}

\begin{softbox}{milcGradeDark}{milcGradeSoft}{Saber ancestral}
En Colombia, la comunidad no espera a que aparezca un testigo: lo nombra. Las juntas de acción comunal de barrios, corregimientos y veredas campesinas postulan a los \textbf{conciliadores en equidad}. Lo que un conciliador hace constar en un acta tiene efecto de cosa juzgada (Ley 446 de 1998). En cada junta hay una \textbf{Comisión de Convivencia y Conciliación}. No se puede sancionar a nadie sin pasar por ella (Ley 2166 de 2021, art. 29). Y los jueces de paz se eligen por voto popular y fallan en equidad (Ley 497 de 1999). Mirar, aquí, es un cargo. Ojo con una cosa: esto no es una tradición antigua, es derecho colombiano vigente. En un barrio de Cartago o en una vereda del Valle, quien da fe fue escogido por sus vecinos. La \textbf{cara de exclusión}: el que da fe también puede callar. Donde el poder local es fuerte, mirar y responder tiene costo. Hoy vas a hacer lo mismo con el ciberbullying: repartir el trabajo de mirar, en vez de dejárselo a la víctima.\par\smallskip{\footnotesize\textcolor{milcNegro!70}{\textbf{Fuente.} Congreso de la República de Colombia. (2021, 18 de diciembre). \emph{Ley 2166 de 2021. Por la cual se establece el régimen de los organismos de acción comunal\ldots} Diario Oficial.}}
\end{softbox}

\begin{softbox}{milcTurquesa}{milcGris}{Saber contemporáneo}
El \textbf{ciberbullying} es el hostigamiento repetido de una persona por medios digitales. Cinco rasgos lo separan de una pelea cualquiera. \textbf{Repetición}: pasa varias veces, no una. \textbf{Desbalance de poder}: quien agrede tiene ventaja, casi siempre por el grupo que lo respalda. \textbf{Intención de daño}: busca hacer sufrir. \textbf{Espacio digital}: ocurre en chats, redes o juegos en línea. \textbf{Permanencia}: lo publicado deja rastro y el daño sigue. Colombia tiene canales abiertos para responder. La \textbf{Línea 141} del ICBF atiende gratis, las 24 horas, a cualquier menor en riesgo. \textbf{Te Protejo} (teprotejo.org) recibe denuncias de contenido, incluso anónimas. La \textbf{Policía} atiende por el 123 lo que necesita presencia inmediata. Y tu colegio debe tener un protocolo escrito, según la \textbf{Ley 1620 de 2013}. Lo importante es esto: los cuatro canales atienden a testigos, no solo a víctimas.
\end{softbox}

\begin{softbox}{milcMostaza}{milcGris}{Pregunta-puente}
\textit{La junta de acción comunal nombra a alguien para que mire y responda, y así nadie carga solo con eso. Cuando en un chat de tu grupo empieza el hostigamiento, ¿quién tiene ese cargo?}
\end{softbox}

\begin{softbox}{milcVerde}{milcGris}{Saber y saber hacer}
Al terminar podrás: (1) \textbf{identificar} si una situación cumple los cinco rasgos del ciberbullying o es un conflicto cotidiano; (2) \textbf{analizar} las cuatro fases y decidir qué canal corresponde en cada una; (3) \textbf{crear} tres piezas que le sirvan a tu grupo: el mapa de las fases, tu plan personal con nombres reales y el cartel del aula.
\end{softbox}



\small
\begin{tabularx}{\textwidth}{>{\bfseries}p{3.0cm}Y}
\rowcolor{milcGradePrimary!12}Autor & Dr. Álvaro Cárdenas Orozco\\
DBA & Producción multimedia ética y comportamiento responsable en entornos digitales (MEN, T\&I)\\
Referentes & Dussel (estética de la liberación) · Ley 1336 · Floridi (infoesfera)\\
Producto final & Tres piezas de convivencia digital: (a) el mapa de las cuatro fases del ciberbullying, de las primeras señales a la denuncia, con dos ejemplos concretos por fase; (b) un plan personal de cinco pasos con nombres reales de las personas a las que acudirías; (c) un cartel para el aula con la Línea 141 destacada y los demás canales, listo para colgar. Cierra con una ficha de cinco líneas donde declaras un compromiso con la convivencia del grupo.\\
\end{tabularx}
\normalsize

\puente{En las sesiones 1 a 5 aprendiste a producir imagen, relato y audio. Hoy esas mismas herramientas se miran desde la convivencia: sirven para sostener a alguien o para hundirlo. La sesión 7 trata grooming y sexting desde la ley; la 8, la estética de la liberación.}

\milcestacion{milcGradeDark}{Ruta MILC: así vamos a aprender}
\begin{tabularx}{\textwidth}{>{\centering\arraybackslash}X>{\centering\arraybackslash}X>{\centering\arraybackslash}X>{\centering\arraybackslash}X}
\phasecard{1}{milcVerde}{milcGris}{Escucha\\[-1mm]\small Observo, escucho y pregunto.} &
\phasecard{2}{milcTurquesa}{milcGris}{Entender\\[-1mm]\small Sistematizo y ordeno.} &
\phasecard{3}{milcMagenta}{milcGris}{Hacer\\[-1mm]\small Construyo y aplico.} &
\phasecard{4}{milcVino}{milcGris}{Cierre\\[-1mm]\small Evalúo y reflexiono.}
\end{tabularx}


\puente{Antes de la teoría vas a clasificar tres situaciones con los cinco rasgos. Ninguna es de tu grupo: son casos hipotéticos que entrega el docente. La idea es que aprendas a distinguir antes de que te toque decidir en caliente.}

\milcestacion{milcVerde}{Estación 1 de 3 · Escucha}

\begin{softbox}{milcVerde}{milcGris}{Escena para escuchar}
\textbf{Actividad 1 · IDENTIFICA --- Las 5 propiedades} (15 min · individual). (1) Tu docente entrega tres situaciones hipotéticas de conflicto digital. (2) Para cada una, marca cuáles de los cinco rasgos cumple: repetición, desbalance de poder, intención de daño, espacio digital y permanencia. (3) Escribe cuál de las tres es ciberbullying en sentido estricto y cuál es un conflicto cotidiano. (4) Anota en una línea qué rasgo hace que el daño sea más difícil de reparar.
\end{softbox}

{\sffamily\bfseries\fontsize{8}{10}\selectfont\textcolor{milcNegro!55}{MARCA CUANDO LO HAYAS HECHO}}
\milccheckrow{Marqué los cinco rasgos en las tres situaciones.}
\milccheckrow{Sé cuál de las tres no es ciberbullying y por qué.}
\milccheckrow{Puedo nombrar el rasgo que hace más difícil reparar el daño.}

\begin{infoband}{milcVerde}


\textbf{Lo que va al cuaderno (Actividad 1 · IDENTIFICA).} \textbf{Título:} «Actividad 1 --- Las 5 propiedades». \textbf{Formato:} tabla de 3 filas y 5 columnas, una fila por situación y una columna por rasgo, con marca donde cumpla. \textbf{Extensión:} un tercio de página. \textbf{Sabes que terminaste cuando} las tres filas están marcadas y escribiste cuál no es ciberbullying.
\end{infoband}

\puente{Ya sabes reconocer. Ahora vas a ver cómo escala el daño en cuatro fases y qué canal corresponde a cada una, con su número o su enlace.}

\milcestacion{milcTurquesa}{Estación 2 de 3 · Entender}

\begin{softbox}{milcTurquesa}{milcGris}{Lo que vas a entender}
\textbf{Actividad 2 · ANALIZA --- Las 4 fases y los canales} (30 min · parejas). (1) Con tu pareja, escriban las cuatro fases en orden: primeras señales, escalada, impacto, denuncia y apoyo. (2) Pongan dos ejemplos concretos en cada fase, escritos como los diría alguien de su edad. (3) Al lado de cada fase, anoten qué canal corresponde y con qué dato de contacto. (4) Marquen con una flecha en qué fase la intervención cuesta menos y explica por qué en una línea.
\end{softbox}

\milcpaso{1}{milcTurquesa}{\textbf{Las cuatro fases.} \textbf{Primeras señales}: un comentario cargado, una broma con filo, un emoji burlón. \textbf{Escalada}: se vuelve repetido y la persona empieza a evitar el chat o el salón. \textbf{Impacto}: aparecen el aislamiento, la irritabilidad, el bajón académico o los problemas de sueño. \textbf{Denuncia y apoyo}: se activa un canal. Cuanto antes se llegue a la cuarta, menos daño hubo.}
\milcpaso{2}{milcTurquesa}{\textbf{Los cuatro canales, con su dato.} \textbf{Línea 141} (ICBF): gratuita, 24 horas, para cualquier menor en riesgo. Contesta un profesional que escucha, orienta y deriva. \textbf{Te Protejo} (teprotejo.org): para denunciar contenido, y acepta denuncia anónima. \textbf{Policía} (123 o el CAI más cercano): cuando hace falta presencia inmediata. \textbf{Protocolo del colegio} (Ley 1620 de 2013): el primer contacto suele ser el docente o el coordinador.}
\milcpaso{3}{milcTurquesa}{\textbf{El testigo también llama.} Ninguno de los cuatro canales exige ser la víctima. Un amigo, un hermano o un compañero de salón puede llamar al 141 y contar lo que vio. Es exactamente lo que hace la Comisión de Convivencia de una junta: recibir de quien mira, no solo de quien sufre.}
\milcpaso{4}{milcTurquesa}{\textbf{La evidencia se guarda antes, no después.} Captura de pantalla, fecha y hora. Sin eso, una denuncia se vuelve la palabra de uno contra la del otro. Guardar no es acusar: es dejar constancia, igual que el acta del conciliador.}

\begin{drawbox}{milcTurquesa}{Ciberbullying y denuncia, en una mirada}
\textbf{5 rasgos:} repetición + poder + intención + digital + permanencia. \\[1mm] \textbf{4 fases:} señales / escalada / impacto / denuncia. \\[1mm] \textbf{4 canales:} 141 / teprotejo.org / 123 / protocolo del colegio. \\[1mm] \textbf{Regla:} el testigo también llama.
\end{drawbox}

\begin{softbox}{milcMostaza}{milcGris}{Errores comunes y su corrección}
(1) \textbf{Llamar bullying a cualquier pelea}: sin los cinco rasgos, es un conflicto y se trata como conflicto. (2) \textbf{Esperar a la fase de impacto}: cuando ya hay aislamiento, el daño está hecho. (3) \textbf{Creer que solo la víctima denuncia}: los cuatro canales reciben a testigos. (4) \textbf{No guardar capturas}: el contenido se borra y la denuncia queda sin sustento. (5) \textbf{Reenviar «para que se vea»}: cada reenvío amplía la audiencia del daño.
\end{softbox}

\begin{infoband}{milcTurquesa}


\textbf{Lo que va al cuaderno (Actividad 2 · ANALIZA).} \textbf{Título:} «Actividad 2 --- Las 4 fases y los canales». \textbf{Formato:} tabla de 4 filas y 3 columnas (fase / dos ejemplos / canal con su dato de contacto), más la flecha en la fase donde intervenir cuesta menos. \textbf{Extensión:} media página. \textbf{Sabes que terminaste cuando} las cuatro filas tienen ejemplos propios y cada canal aparece con su número o su enlace.
\end{infoband}

\puente{Tienes el criterio y los canales. Toca construir: el mapa de las fases, tu plan personal con nombres reales y el cartel que se cuelga mañana en el aula.}

\milcestacion{milcMagenta}{Estación 3 de 3 · Hacer}

\begin{softbox}{milcMagenta}{milcGris}{Cómo vas a construir tu producto}
\textbf{Actividad 3 · CREA --- Cartel, plan y compromiso} (25 min · individual). (1) Pasa el mapa de las cuatro fases a limpio, en papel o en una herramienta libre. (2) Escribe tu plan personal de cinco pasos con nombres reales de personas de tu vida. (3) Diseña el cartel del aula con la jerarquía visual de la sesión 1, con la Línea 141 como el elemento más grande. (4) Escribe tu compromiso en cinco líneas. (5) Pásale el cartel a un compañero y pregúntale a quién llamaría; anota si acertó.
\end{softbox}

\milcpaso{1}{milcMagenta}{\textbf{Tu entrega tiene cuatro piezas.} El mapa de las cuatro fases con dos ejemplos cada una. El plan personal de cinco pasos. El cartel del aula con los cuatro canales. La ficha de compromiso de cinco líneas.}
\milcpaso{2}{milcMagenta}{\textbf{Nombres reales, no categorías.} «Hablo con alguien de confianza» no es un plan. «Hablo con mi tía Lucía» sí lo es. Decidir hoy a quién llamarías mañana es la mitad del trabajo, porque en el momento de la crisis ya no se decide bien.}
\milcpaso{3}{milcMagenta}{\textbf{El cartel se juzga desde lejos.} Párate a tres metros. Si desde ahí no se lee el 141, la jerarquía está mal. Lo primero que se ve tiene que ser el número, no el título ni el dibujo.}
\milcpaso{4}{milcMagenta}{\textbf{El compromiso se escribe en primera persona y con verbo.} «Me comprometo a no reenviar capturas burlonas de compañeros.» «Me comprometo a preguntarle a quien vea aislado.» Una promesa sin verbo no se puede cumplir ni revisar.}

\begin{drawbox}{milcMagenta}{Antes de entregar}
\checkbox Mapa con las cuatro fases y dos ejemplos en cada una. \\[1mm] \checkbox Plan personal de cinco pasos con nombres reales. \\[1mm] \checkbox Cartel con la Línea 141 como elemento más grande. \\[1mm] \checkbox Los cuatro canales aparecen en el cartel con su dato. \\[1mm] \checkbox Compromiso de cinco líneas, en primera persona y con verbo. \\[1mm] \checkbox Un compañero miró el cartel y supo a quién llamar.
\end{drawbox}

\begin{softbox}{milcMagenta}{milcGris}{Plantilla de trabajo}
\textbf{Mapa.} \textit{Cuatro fases, dos ejemplos por fase.} \\[1mm] \textbf{Plan.} \textit{Cinco pasos con nombres reales.} \\[1mm] \textbf{Cartel.} \textit{141 grande, más los otros tres canales.} \\[1mm] \textbf{Compromiso.} \textit{Cinco líneas en primera persona.}
\end{softbox}

\begin{infoband}{milcMagenta}


\textbf{Lo que va al cuaderno (Actividad 3 · CREA).} \textbf{Título:} «Actividad 3 --- Cartel, plan y compromiso». \textbf{Formato:} el plan de cinco pasos con nombres, el boceto del cartel y las cinco líneas del compromiso. \textbf{Extensión:} media página. \textbf{Sabes que terminaste cuando} un compañero miró tu cartel y dijo a quién llamaría sin preguntarte nada.
\end{infoband}

\puente{Lo que entregas hoy: mapa, plan, cartel y compromiso. La prueba de calidad es sencilla: un compañero que no estuvo en clase mira el cartel y sabe a quién llamar.}

\milcestacion{milcMostaza}{Producto de la sesión}
\textbf{Mapa de las 4 fases + plan personal con nombres + cartel con la Línea 141 + compromiso}.
\begin{softbox}{milcMostaza}{milcGris}{Criterios de calidad}
(1) El mapa tiene las cuatro fases y dos ejemplos propios en cada una. (2) El plan de cinco pasos nombra personas reales, no categorías. (3) El cartel se lee a tres metros y la Línea 141 es lo primero que se ve. (4) Los cuatro canales aparecen con su número o su enlace. (5) El compromiso son cinco líneas en primera persona y con verbo.
\end{softbox}

\puente{Antes de cerrar, mira lo que hiciste desde cinco lados. Conocer un canal de denuncia y no decirlo es dejar sola a la persona que lo necesitaba.}

\milcestacion{milcVino}{Evaluación liberadora · cinco dimensiones}

\begin{softbox}{milcVino}{milcGris}{Sentido de la evaluación}
Evaluar no es solo revisar una nota. Es reconocer qué aprendí, qué cambió en mí, qué emociones manejé, cómo me relaciono con la ciudad y la comunidad, y qué le dejaré a quien viene detrás.
\end{softbox}

\textbf{Mi reflexión liberadora en cinco dimensiones.} Completa cada frase iniciada:

\small
\begin{tabularx}{\textwidth}{>{\bfseries\RaggedRight\arraybackslash}p{3.4cm}Y>{\RaggedRight\arraybackslash}p{4.6cm}}
\rowcolor{milcVino}\color{white}Dimensión & \color{white}\bfseries Mi respuesta (frame starter) & \color{white}\bfseries Algo que descubrí\\
1. Personal & Lo que más cambió en mí fue \linead{6.5cm} & \linead{4.2cm}\\
2. Emocional & La emoción más fuerte fue \linead{2.5cm} y la regulé \linead{3cm} & \linead{4.2cm}\\
3. Ciudadana & En democracia esto importa porque \linead{6cm} & \linead{4.2cm}\\
4. Local (Cartago) & En mi barrio sería útil para \linead{6.5cm} & \linead{4.2cm}\\
5. Intergeneracional & A mi abuela/abuelo le diría \linead{6.5cm} & \linead{4.2cm}\\
\end{tabularx}
\normalsize

\begin{infoband}{milcVino}
\textbf{Para la próxima sustentación.} Vuelve a esta tabla en seis meses. Verás que las respuestas cambian — no porque lo aprendido cambie, sino porque tú cambias. La evaluación liberadora es una conversación contigo a través del tiempo.
\end{infoband}


\puente{Para terminar, tres ideas sobre la responsabilidad de mirar, contadas con palabras de esta guía: Dussel sobre escuchar la voz del que reclama, Marco Aurelio sobre el enjambre y la abeja, Floridi sobre vivir entre otros en un mismo entorno de información.}

\milcestacion{milcGradeDark}{Tres ideas para cerrar · Dussel, un estoico y Floridi}

\begin{softbox}{milcVino}{milcGris}{Enrique Dussel · lente del nosotros}
\textit{Tener conciencia ética es ser capaz de oír la voz del que reclama, aunque incomode a las reglas del grupo.}

{\footnotesize\itshape\color{milcNegro!70}--- idea tomada de Enrique Dussel, contada con palabras de esta guía. No es una cita textual.}

En un chat, la persona hostigada casi nunca grita. Dice «ya, déjenme» y el grupo sigue. Oír esa voz cuando todos la están tapando es lo que hoy entrenaste con el plan personal.

\textbf{¿Cuándo oí a alguien pedir que pararan y seguí como si nada?}
\end{softbox}
\linead{\linewidth}\\[1mm]
\linead{\linewidth}

\begin{softbox}{milcMostaza}{milcGris}{Estoicismo · Marco Aurelio · cuidado interior}
\textit{Lo que no le sirve al enjambre tampoco le sirve a la abeja.}

{\footnotesize\itshape\color{milcNegro!70}--- idea tomada de Marco Aurelio, contada con palabras de esta guía. No es una cita textual.}

Un grupo donde se puede hostigar a cualquiera no es seguro para nadie, ni siquiera para quien se ríe hoy. El cartel que cuelgas no protege solo a una persona: cambia lo que el salón acepta.

\textbf{¿Lo que dejo pasar en mi grupo me protegería a mí si mañana me tocara?}
\end{softbox}
\linead{\linewidth}\\[1mm]
\linead{\linewidth}

\begin{softbox}{milcTurquesa}{milcGris}{Luciano Floridi · lente de la infoesfera}
\textit{Vivimos dentro de un mismo entorno de información, y compartirlo con otros nos hace responsables de lo que dejamos allí.}

{\footnotesize\itshape\color{milcNegro!70}--- idea tomada de Luciano Floridi, contada con palabras de esta guía. No es una cita textual.}

Una captura burlona no desaparece cuando la borras: alguien ya la reenvió. Por eso la permanencia es el rasgo más difícil de reparar, y por eso no reenviar es una decisión ética, no un gesto menor.

\textbf{¿Qué dejé en ese entorno esta semana que no me gustaría encontrar sobre mí?}
\end{softbox}
\linead{\linewidth}\\[1mm]
\linead{\linewidth}

\begin{infoband}{milcNegro}
\textbf{Nota para el docente.} Las tres ideas están contadas con palabras de esta guía a partir del pensamiento de cada autor; no son citas textuales. De 9.º en adelante se trabajan con la obra y su referencia.
\end{infoband}

\milcestacion{milcVino}{Mi autoevaluación · marca una casilla por fila}

{\small
\setlength{\extrarowheight}{2.2pt}
\begin{tabularx}{\textwidth}{>{\RaggedRight\arraybackslash\bfseries}p{3.0cm}YYY}
\rowcolor{milcVino}\color{white}\bfseries Criterio & \color{white}\bfseries Lo logré & \color{white}\bfseries En proceso & \color{white}\bfseries Necesito apoyo\\[.6mm]
Comprensión del tema & \milcopcion{Explico las ideas con mis palabras.} & \milcopcion{Reconozco las ideas, falta precisión.} & \milcopcion{Necesito repasar lo básico.}\\[.8mm]
\rowcolor{milcGris}Pensamiento computacional & \milcopcion{Usé los pilares al armar mi producto.} & \milcopcion{Usé algunos pilares.} & \milcopcion{Necesito releer la estación 2.}\\[.8mm]
Aplicación y producto & \milcopcion{Mi producto cumple los criterios.} & \milcopcion{Está armado, con ajustes pendientes.} & \milcopcion{Aún está en construcción.}\\[.8mm]
\rowcolor{milcGris}Reflexión 5 dimensiones & \milcopcion{Respondí cada una con honestidad.} & \milcopcion{Respondí algunas con detalle.} & \milcopcion{Mis respuestas son muy generales.}\\[.8mm]
Triángulo de pensamiento & \milcopcion{Conecté las 3 voces con mi trabajo.} & \milcopcion{Conecté al menos una.} & \milcopcion{No las relaciono con mi proceso.}\\[.8mm]
\end{tabularx}}

\vspace{3mm}
\textbf{Hoy entendí que:}\\[1mm]
\linead{\linewidth}\\[1mm]
\linead{\linewidth}

\textbf{Mi compromiso para la próxima sesión es:}\\[1mm]
\linead{\linewidth}\\[1mm]
\linead{\linewidth}

\begin{softbox}{milcMostaza}{milcGris}{Cita de despedida · Marco Aurelio}
\textit{``El obstáculo en el camino se vuelve el camino. Lo que estorba al camino se vuelve camino.''}\\[1mm]
Cada error de hoy, bien observado, se convierte en pieza del camino que pisarás mañana.
\end{softbox}

\small
\begin{tabularx}{\textwidth}{>{\bfseries}p{3.6cm}Y}
\rowcolor{milcGradeDark!12}Nombre completo: & \linead{10cm}\\
Fecha: & \linead{4cm} \quad Curso/grupo: \linead{4cm}\\
Mi firma: & \linead{9cm}\\
\end{tabularx}
\normalsize

\begin{infoband}{milcGradeDark}
\textbf{Para la próxima sesión.} Llega con el mapa, el plan, el cartel y las cinco líneas del compromiso. El cartel se cuelga en el aula al empezar la clase. En la sesión 7 vas a trabajar grooming y sexting desde la ley y los protocolos de prevención. Es la continuación directa de lo de hoy.
\end{infoband}

\end{document}
